% In this file you should put the actual content of the blueprint.
% It will be used both by the web and the print version.
% It should *not* include the \begin{document}

\chapter{The Revelation Principle for communication equilibria}

Players $i \in N$ have finite type sets $T_i$ and action sets $A_i$. A
communication device has finite message sets $M_i$ and signal sets $R_i$.
Following Myerson (1982) and Forges (1986), we show that every
communication-equilibrium outcome can be achieved by truthful-obedient play
under a direct device.

\section{Finite distributions}

\begin{definition}[Finite distribution]
  \label{def:findist}
  \lean{FinDist}
  \leanok
  A distribution $\mu$ on a finite set $X$ is a weight $\mu(x) \ge 0$ for each
  $x \in X$, with $\sum_x \mu(x) = 1$.
\end{definition}

\begin{definition}[Point mass, mixture, pushforward]
  \label{def:findist-ops}
  \lean{FinDist.pure, FinDist.bind, FinDist.map}
  \leanok
  \uses{def:findist}
  The point mass $\delta_a$; the mixture
  $(\mu \gg\!= f)(b) = \sum_a \mu(a)\, f(a)(b)$; and the pushforward
  $f_*\mu$.
\end{definition}

\begin{definition}[Expectation]
  \label{def:findist-E}
  \lean{FinDist.E}
  \leanok
  \uses{def:findist}
  $\mathbb{E}_\mu[g] = \sum_x \mu(x)\, g(x)$.
\end{definition}

\begin{definition}[Independent products]
  \label{def:findist-prod}
  \lean{FinDist.pi, FinDist.prod}
  \leanok
  \uses{def:findist}
  The product of a family, $(\otimes_i \mu_i)(s) = \prod_i \mu_i(s_i)$, and of
  a pair, $(\mu \otimes \nu)(a,b) = \mu(a)\,\nu(b)$.
\end{definition}

\section{Games, devices and equilibrium}

\begin{definition}[Bayesian game]
  \label{def:bayes-game}
  \lean{BayesGame}
  \leanok
  \uses{def:findist}
  A Bayesian game $G$ consists of a common prior $p$ on type profiles and
  payoffs $u_i(a, t)$.
\end{definition}

\begin{definition}[Communication device]
  \label{def:device}
  \lean{Device}
  \leanok
  \uses{def:findist}
  A device $d$ maps each message profile $m$ to a distribution $q(\cdot \mid m)$
  of signal profiles.
\end{definition}

\begin{definition}[Strategies]
  \label{def:strategies}
  \lean{ReportRule, ActionRule, PureStrat, MixStrategyProfile}
  \leanok
  \uses{def:findist}
  A pure strategy of player $i$ is a report rule $T_i \to M_i$ and an action
  rule $T_i \times M_i \times R_i \to A_i$. A mixed strategy profile assigns
  each player a distribution over pure strategies.
\end{definition}

\begin{definition}[Play and outcome]
  \label{def:outcome}
  \lean{play, outcome}
  \leanok
  \uses{def:device, def:strategies, def:findist-ops, def:findist-prod}
  Given types $t$, players draw pure strategies independently, send messages,
  receive signals from $d$, and act. The outcome $\mu(\cdot \mid t)$ is the
  resulting distribution of action profiles.
\end{definition}

\begin{definition}[Payoff]
  \label{def:payoff}
  \lean{payoff}
  \leanok
  \uses{def:bayes-game, def:outcome, def:findist-E}
  Player $i$'s ex-ante payoff is
  $\mathbb{E}_{t \sim p}\, \mathbb{E}_{a \sim \mu(\cdot \mid t)}[u_i(a, t)]$.
\end{definition}

\begin{definition}[Communication equilibrium]
  \label{def:comm-eq}
  \lean{IsCommunicationEq}
  \leanok
  \uses{def:payoff}
  $\pi$ is a communication equilibrium of $G(d)$ if no player $i$ gains by
  switching to another mixed strategy $\rho$.
\end{definition}

\begin{definition}[Communication-equilibrium outcome]
  \label{def:comm-eq-outcome}
  \lean{IsCommunicationEqOutcome}
  \leanok
  \uses{def:comm-eq}
  $\mu$ is a communication-equilibrium outcome if some device $d$ and
  communication equilibrium $\pi$ of $G(d)$ produce it.
\end{definition}

\begin{definition}[Direct device, truthful play]
  \label{def:direct}
  \lean{DirectDevice, truthful, truthfulProfile}
  \leanok
  \uses{def:device, def:strategies}
  A direct device has $M_i = T_i$ and $R_i = A_i$. Truthful play reports the
  true type and obeys the recommended action.
\end{definition}

\section{The proof}

\begin{definition}[Canonical device]
  \label{def:canonical}
  \lean{canonical}
  \leanok
  \uses{def:direct, def:outcome}
  Given $d$ and $\pi$, the canonical device $d^*$ takes reported types $t$ and
  recommends actions drawn from $\mu(\cdot \mid t)$.
\end{definition}

\begin{definition}[Mimicking strategy]
  \label{def:deviate}
  \lean{deviate, deviateMixed}
  \leanok
  \uses{def:strategies, def:findist-prod}
  Given a strategy $x$ of player $i$ under $d^*$ and a strategy $y$ under $d$,
  the mimic plays $y$ as if the type were $x$'s report, then applies $x$'s
  action rule to $y$'s action.
\end{definition}

\begin{lemma}[Truthful outcome equivalence]
  \label{lem:truthful}
  \lean{canonical_truthful_outcome_equivalence}
  \leanok
  \uses{def:canonical}
  Truthful play under $d^*$ has the same outcome as $\pi$ under $d$.
\end{lemma}
\begin{proof}
  The truthful profile is a point mass, so reports equal $t$; then $d^*$
  draws from $\mu(\cdot \mid t)$, and obeying is the identity.
\end{proof}

\begin{lemma}[Deviation outcome equivalence]
  \label{lem:deviation}
  \lean{canonical_deviation_outcome_equivalence}
  \leanok
  \uses{def:canonical, def:deviate}
  Deviating to $\rho$ under $d^*$ has the same outcome as playing the mimic
  of $\rho$ under $d$.
\end{lemma}
\begin{proof}
  Both sides are nested independent draws of $x \sim \rho$ and
  $y \sim \otimes_j \pi_j$; swap their order, and note that redrawing player
  $i$'s coordinate from $\pi_i$ does not change the distribution.
\end{proof}

\begin{theorem}[Revelation Principle, canonical form]
  \label{thm:revelation-canonical}
  \lean{revelation_principle_canonical}
  \leanok
  \uses{def:comm-eq, def:canonical}
  If $\pi$ is a communication equilibrium of $G(d)$, then truthful play is a
  communication equilibrium of $G(d^*)$ with the same outcome.
\end{theorem}
\begin{proof}
  \leanok
  \uses{lem:truthful, lem:deviation}
  A deviation's payoff under $d^*$ equals its mimic's payoff under $d$, which
  is at most the equilibrium payoff, which equals truthful play's payoff.
\end{proof}

\begin{theorem}[Revelation Principle, existence form]
  \label{thm:revelation}
  \lean{revelation_principle}
  \leanok
  \uses{def:comm-eq, def:direct}
  If $\pi$ is a communication equilibrium of $G(d)$, some direct device has
  truthful play as a communication equilibrium with the same outcome.
\end{theorem}
\begin{proof}
  \leanok
  \uses{thm:revelation-canonical}
  Take the canonical device.
\end{proof}

\begin{theorem}[Revelation Principle]
  \label{thm:revelation-iff}
  \lean{revelation_principle_iff}
  \leanok
  \uses{def:comm-eq-outcome, def:direct}
  $\mu$ is a communication-equilibrium outcome if and only if $\mu$ is the
  outcome of a truthful equilibrium under some direct device.
\end{theorem}
\begin{proof}
  \leanok
  \uses{thm:revelation}
  ($\Rightarrow$) by the existence form; ($\Leftarrow$) a direct device is a
  device.
\end{proof}
